% !TEX root = ../bb_AiRa_Kappaleet.tex

% ö = \%c3\%b6
% ä = \%C3\%A4

\xdef\sectiontitle{\Isectiontitle} %aiheuttama
\TOCrefs

%%%%%%%%%%%%%%%
\begin{frame}[label=1_8_a]%[label=1_8_TOC1]
\frametitle{\texorpdfstring{\culine[\sectiontitleulinecolor]{\sectiontitle}}{\sectiontitle} }
\label{\TOCname}

\vspace{-0.5cm}

\begin{itemize}[leftmargin=0.5cm,itemsep=2mm]
    \item[1.1]<1-> \hyperlink{1_1_a}{\IaSubsectiontitle}
    \item[1.2]<1-> \hyperlink{1_2_a}{\IbSubsectiontitle}	
    \item[1.3]<1-> \hyperlink{1_3_a}{\IcSubsectiontitle}	
    \item[1.4]<1-> \hyperlink{1_4_a}{\IdSubsectiontitle}
    \item[1.5]<1-> \hyperlink{1_5_a}{\IeSubsectiontitle}
    \item[1.6]<1-> \hyperlink{1_6_a}{\IfSubsectiontitle}
    \item[1.7]<1-> \hyperlink{1_7_a}{\IgSubsectiontitle}
\end{itemize}

\vspace{-1cm}
\begin{itemize}[leftmargin=8.5cm,itemsep=2mm]
    \item[1.8]<1-> \hyperlink{1_8_a}{\CDAlert[red]{\IhSubsectiontitle}}
	\item[1.9]<1-> \hyperlink{1_9_a}{\IiSubsectiontitle}
	\item[1.10]<1-> \hyperlink{1_10_a}{\IjSubsectiontitle}
	\item[1.11]<1-> \hyperlink{1_11_a}{\IkSubsectiontitle}
\end{itemize}

%\endofslide 
\end{frame}
%%%%%%%%%%%%%%%


\xdef\sectiontitle{\IhSubsectiontitle} 

%%%%%%%%%%%%%%%
\begin{frame}
\frametitle{\texorpdfstring{\culine[\sectiontitleulinecolor]{\sectiontitle}}{\sectiontitle} }
\framesubtitle{Tilatiheys}
\appear{}
\begin{itemize}
	\item Toisin kuin elektronit (ja He-4 atomit), fotonit ovat bosoneita (spin $=1$) 
	%\item The behavior of  
	\item Fotonikaasun kuvaamiseksi\appear{ meidän täytyy laskea fotonien \CDAlert[fuchsia]{tilatiheys}}

\item Fotonien energia liittyy suoraan niiden taajuuteen $f$ \appear{(eikä niinkään energiatason kvanttilukuun $n$):} $ \qquad \appear{E =h f } \equal \appear{h \frac{c}{\lambda}} $
\end{itemize}

\newcommand\myeq{\stackrel{\mathclap{\footnotesize\mbox{suhtautuu}}}{~\Leftrightarrow~}}
\begin{minipage}{8cm}
\begin{itemize}

	\item Tarkastellaan aluksi 3D-kvanttikaivoa jonka leveys on  $L$\appear{, mihin mahtuu pie\-nim\-millään $\lambda=2 L$ kokoinen kvantti/fotoni}
	\appear{\xdef\loc{south east}%
\begin{tikzpicture}[remember picture,overlay] % anchor=south
    \node (A) [yshift=-0.25*\figYshift,xshift=\figXshift, anchor=\loc] at (current page.\loc) {\includegraphics[page=1,width=5.2cm]{Figures_booklet/SOM_Tipler_Statistics_and_Bonding_figures/SOM_Tipler_figure.png}};%scale=\figscale. % 8cm max hei
\end{tikzpicture}}%
	\[
\begin{aligned}
n_{x}^{2} \plus \appear{n_{y}^{2}} \plus \appear{n_{z}^{2}}  \equal  \appear{n^{2}}  \qquad \myeq \qquad  \appear{R^{2}}  \equal  \appear{ \bigg( \Frac{L}{\lambda / 2}  \bigg)^{2}}
\end{aligned}
\]
\item[] jonka avulla tilojen $n$ täyttämän pallon säde tila-avaruudessa  voidaan laskea:
\[
R  \equal  \frac{L}{\lambda / 2}  \equal  \frac{2 L}{\lambda}  \equal  \frac{2 L E}{h c}
\]
\end{itemize}
\end{minipage}

\endofslide 
%\endofsection
\end{frame}
%%%%%%%%%%%%%%%

%%%%%%%%%%%%%%%
\begin{frame}
\frametitle{\texorpdfstring{\culine[\sectiontitleulinecolor]{\sectiontitle}}{\sectiontitle} }
\framesubtitle{Tilatiheys}

\appear{}
\begin{itemize}
	\item Montako 'hilapistettä' \appear{mahtuu $n$-avaruuden palloon jonka säde on  $R$?}
\[
\begin{aligned}
& & R& \equal \frac{2 L E}{h c} \\ 
&\appear{\Rightarrow} \qquad& \appear{N} &\equal \appear{2 \cdot} \frac{1}{8} \appear{\cdot} \frac{4}{3} \appear{\pi R^{3}}  \equal  \frac{\pi}{3} \appear{ \bigg( \frac{2 L E}{h c}} \appear{ \bigg)^{3}}  
\equal  \frac{8 \pi}{3} \frac{E^{3}}{h^{3} c^{3}} \appear{L^{3}}  
\equal \frac{8 \pi}{3} \frac{E^{3}}{h^{3} c^{3}} \appear{V}
\end{aligned}
\]

\item[] missä $N$ on tilojen lukumäärä \appear{säteeltään $R=2 L E /(h c)$ olevan pallon sisällä}
\item Nyt energiavälillä $d E$ olevien pisteiden lukumäärä voidaan laskea:
\[
\frac{d N}{d E}  \equal \frac{8 \pi}{3} \frac{3 E^{2}}{h^{3} c^{3}} \appear{V}
 \equal \frac{8 \pi E^{2}}{h^{3} c^{3}} \appear{V} 
 \qquad \Rightarrow \qquad 
 \frac{d N}{d f}  \equal \frac{8 \pi f^{2}}{c^{3}} \appear{V} \qquad \appear{(E=h f} \quad \Rightarrow \quad \appear{d E=h d f)}
\]

\end{itemize}

\endofslide 
%\endofsection
\end{frame}
%%%%%%%%%%%%%%%

%%%%%%%%%%%%%%%
\begin{frame}
\frametitle{\texorpdfstring{\culine[\sectiontitleulinecolor]{\sectiontitle}}{\sectiontitle} }
\framesubtitle{Tilatiheys}

\appear{}
\begin{itemize}[itemsep=2.5mm]
	\item Eli olemme laskeneet tilojen lukumäärän \appear{jotka ovat tietyllä taajuusvälillä:}
\[
\frac{\text { differentiaalinen tilojen lukumäärä taajuusvälillä }}{d f}  \equal \frac{8 \pi V}{c^{3}} \appear{f^{2}}
\]
\item Tai vastaava energian avulla ilmaistuna:
\[
\frac{\text { differentiaalinen tilojen lukumäärä energiavälillä dE }}{d E / h}  \equal \frac{8 \pi V}{c^{3}} \appear{(E / h)^{2}}
\]

\item Eli fotonikaasun \CDAlert[red]{tilatiheydeksi $D=D(E)$} \appear{(tilavuudessa $V$)} \appear{saadaan:}
\[
D(E)  \equal \frac{8 \pi V}{h^{3} c^{3}} \appear{E^{2}}
\]
\end{itemize}

\endofslide 
%\endofsection
\end{frame}
%%%%%%%%%%%%%%%

%%%%%%%%%%%%%%%
\begin{frame}
\frametitle{\texorpdfstring{\culine[\sectiontitleulinecolor]{\sectiontitle}}{\sectiontitle} }
\framesubtitle{Miehitysluvun normeeraus}
\appear{}
\begin{itemize}[itemsep=1.7mm]
	\item Fotonikaasu eroaa harmonisen värähtelijän tapauksesta\appear{, tai esim. nestemäisen heliumin tapauksesta}

\item Fotonikaasussa \CDAlert[blue]{fotonien lukumäärää ei ole etukäteen rajoitettu}\appear{, systeemi absoboi ja emittoi fotoneita koko ajan}\appear{, fotonien hetkellinen lukumäärä määräytyy systeemin tilan} \appear{ja sen ominaisuuksien mukaan}

\item \CDAlert[red]{Bosen–Einsteinin statistiikka} bosoneille oli muotoa: 
\[ \mathbb{N} (E)_{B E}  \equal  \frac{1}{B e^{E / k_{B} T}-1}
\] 
\item[] missä parametri $B$ täytyisi liittää hiukkasten lukumäärään $N$ 

\item Koska fotonikaasussa $N$ on vapaa muuttuja\appear{, ratkaisu on määritellä $B\equiv\appear{1}$} \appear{(fotonikaasulle)}
\[
 \mathbb{N} (E)_{B E, \text {fotonikaasu}}=\frac{1}{e^{E / k_{B} T}-1}
\]
\end{itemize}

\endofslide 
%\endofsection
\end{frame}
%%%%%%%%%%%%%%%

%%%%%%%%%%%%%%%
\begin{frame}
\frametitle{\texorpdfstring{\culine[\sectiontitleulinecolor]{\sectiontitle}}{\sectiontitle} }
\framesubtitle{Esimerkki}

\appear{}
\begin{itemize}[itemsep=3mm]
	\item \textbf{K}: Vertaa kvantitatiivisesti sähkömagneettisen säteilyn ja aineen sisäisiä energiatiloja
\item \textbf{A}: Voisimme laskea keskimääräiset energiat: 
\[
<E>\,  \equal  \frac{\nint E  \mathbb{N} (E) D(E) d E}{\nint  \mathbb{N} (E) D(E) d E}
\]
\item[] missä osoittaja vastaa systeemin energiaa \appear{(ja nimittäjä vastaa hiukkasten lukumäärää)} 

\item Voisimme myös vain käyttää osoittajaa SMG-säteilyn \appear{(eli fotonikaasun)} \appear{sisäisen energian laskemiseksi} 
\end{itemize}

\endofslide 
%\endofsection
\end{frame}
%%%%%%%%%%%%%%%

%%%%%%%%%%%%%%%
\begin{frame}
\frametitle{\texorpdfstring{\culine[\sectiontitleulinecolor]{\sectiontitle}}{\sectiontitle} }
\framesubtitle{Esimerkki}

\appear{}
\begin{itemize}
%\item But, we could as well use only the nominator to obtain the internal energy of EM-radiation\appear{, or 'photon gas'} 

\item Käyttämällä muuttujanvaihdosta $\appear{E}  \equal \appear{k_{B} T x} \quad \Rightarrow \quad \appear{x}  \equal \appear{E/(k_{B} T)}$\appear{, osoittajaksi saadaan}
\[
\begin{aligned}
U_{fotoni}  &\equal  \appear{\nint_{0}^{\infty}} \appear{E  \mathbb{N} (E)} \appear{D(E)} \appear{d E}
 \equal  
 \appear{\nint_{0}^{\infty} E} \frac{1}{e^{E / k_{B} T}-1} \frac{8 \pi V}{h^{3} c^{3}} \appear{E^{2}} \appear{d E} \\
& \equal 
\frac{8 \pi V}{h^{3} c^{3}} \appear{\nint_{0}^{\infty}} \appear{k_{B} T x} \frac{1}{e^{x}-1} \appear{\textrm{((} k_{B} T x \textrm{))} ^{2}} \appear{k_{B} T} \appear{d x}
 \equal 
 \frac{8 \pi V k_{B}^{4} T^{4}}{h^{3} c^{3}} \appear{\nint_{0}^{\infty}} \frac{x^{3} d x}{e^{x}-1}  \\
 &\equal 
 \frac{8 \pi V k_{B}^{4} T^{4}}{h^{3} c^{3}} \appear{\cdot} \frac{\pi^{4}}{15}
\end{aligned}
\]
\end{itemize}


\endofslide 
%\endofsection
\end{frame}
%%%%%%%%%%%%%%%

%%%%%%%%%%%%%%%
\begin{frame}
\frametitle{\texorpdfstring{\culine[\sectiontitleulinecolor]{\sectiontitle}}{\sectiontitle} }
\framesubtitle{Esimerkki}
\appear{}

\begin{itemize}
	\item Fotonikaasulle saadaan siis
	\[
	U_{fotoni}=\frac{8 \pi V k_{B}^{4} T^{4}}{h^{3} c^{3}} \appear{\cdot} \frac{\pi^{4}}{15}  \equal \frac{8 \pi^{5} V k_{B}^{4} T^{4}}{15 h^{3} c^{3}}
	\]
\item[] $300 \mathrm{~K}$ lämpötilassa \appear{arvoksi saadaan:} $\quad \frac{U_{fotoni}}{V} \appear{=6,1 \times 10^{-6} \mathrm{~J} / \mathrm{m}^{3}}$

\item Tavallisille kaasuille \appear{(\CDAlert[fuchsia]{yksiatominen} ideaalikaasu)} \appear{sisäenergiaksi NTP:ssä saadaan}
\[
\frac{U_{\text {aine}}}{V}
 \equal 
 \frac{N \times \Frac{3}{2} k_{B} T}{V}
  \equal 
  \frac{\Frac{3}{2} p V}{V} \equal \frac{3}{2} \appear{p} \equal \appear{1,5 \times 10^{5} \mathrm{~J} / \mathrm{m}^{3}}
\]
\item Eli suhde on pieni: $\frac{U_{\text {fotoni}}}{U_{\text {aine}}}\appear{=4 \times 10^{-11}$}

 \implies{ Eli NTP-olosuhteissa suhde on sellainen että
 		\Alert[red]{systeemin termo-\\
		dynaamisesta kokonaisenergiasta \appear{erittäin pieni osa on \\
		} \appear{SMG-säteilyä}} }
\end{itemize}

\endofslide 
%\endofsection
\end{frame}
%%%%%%%%%%%%%%%

%%%%%%%%%%%%%%%
\begin{frame}
\frametitle{\texorpdfstring{\culine[\sectiontitleulinecolor]{\sectiontitle}}{\sectiontitle} }
\framesubtitle{Planckin laki, musta kappale ja Wienin siirtymälaki}
\appear{}
%\framesubtitle{Planckin laki, mustan kappaleen säteily ja Wienin laki}

\begin{itemize}
	\item Edellisen esimerkin sisäenergiaa kuvaavan integraalin integrandi voidaan ajatella \appear{systeemin kokonaisenergiana tietyllä fotonienergian välillä} \appear{[[$E , E+d E$]]:}
\[
d U_{fotoni}(E, E+d E)
 \equal 
 \appear{E} \frac{1}{e^{E / k_{B} T}-1} \frac{8 \pi V}{h^{3} c^{3}} \appear{E^{2}} \appear{d E}
  \equal \frac{h f^{3}}{e^{h f / k_{B} T}-1} \frac{8 \pi V}{c^{3}} \appear{d f}
\]
\item[] tämä on \CDAlert[blue]{Planckin laki},  joka antaa meille
\vspace{0.3mm}
\end{itemize}
\begin{minipage}{6.5cm}
\begin{itemize}
	\item[] spektrisen energiatiheyden sähkömagneettiselle säteilylle, tai niinsanotun \CDAlert[red]{mustan kappaleen säteilyn}%
	\appear{\xdef\loc{south east}%
\begin{tikzpicture}[remember picture,overlay] % anchor=south
    \node (A) [yshift=-0.25*\figYshift,xshift=\figXshift, anchor=\loc] at (current page.\loc) {\includegraphics[page=1,width=7.2cm]{Figures_booklet/SOM_9_Statistical_mechanics_figures/SOM_Figure_20.png}}; %scale=\figscale. % 8cm max hei
\end{tikzpicture}%
}
\item Käyrän maksimikohta riippuu lineaarisesti lämpötilasta, eli toteuttaa \CDAlert[tunired]{Wienin siirtymälain} 
\end{itemize}

\end{minipage}


\endofslide 
%\endofsection
\end{frame}
%%%%%%%%%%%%%%%

%%%%%%%%%%%%%%%
\begin{frame}
\frametitle{\texorpdfstring{\culine[\sectiontitleulinecolor]{\sectiontitle}}{\sectiontitle} }
\framesubtitle{Stefanin–Boltzmannin laki}
\appear{}

\begin{itemize}
	\item Käyrän alle jäävän pinta-alan integrointi antaa meille \CDAlert[red]{Stefanin–Boltzmannin lain}:
\[
U \propto \appear{T^{4}}
\]
\item Yhteenvetona\appear{, \CDAlert[fuchsia]{bosoninen fotonikaasun malli}} \appear{mahdollistaa usean erittäin perustavanlaatuisen yhtälön}\appear{/lain johtamisen} \appear{yhdistäen termodynamiikan} \appear{sähködynamiikkaan}
\item[] \begin{itemize}[itemsep=2.5mm]
	 \item Planckin laki
	 \item Mustan kappaleen säteily
	 \item Wienin laki
	 \item Stefanin–Boltzmannin laki
\end{itemize}

\end{itemize}

\xdef\loc{south east}
\begin{tikzpicture}[remember picture,overlay] % anchor=south
    \node (A) [yshift=-0.25*\figYshift,xshift=\figXshift, anchor=\loc] at (current page.\loc) {\includegraphics[page=1,width=7.2cm]{Figures_booklet/SOM_9_Statistical_mechanics_figures/SOM_Figure_20.png}}; %scale=\figscale. % 8cm max hei
\end{tikzpicture}


%\endofslide 
\endofsection
\end{frame}
%%%%%%%%%%%%%%%
